%! AP Statistics — Unit 2, Lesson 2.7 — Student Packet Cover Sheet
\documentclass[10pt]{article}
\usepackage{apstats-article}
\usepackage{apstats-boxes}
\usepackage{ltablex}
\keepXColumns

\begin{document}

% Full-bleed navy banner
\begin{tikzpicture}[remember picture, overlay]
  \fill[navy] (current page.north west)
    rectangle ([yshift=-0.9in]current page.north east);
\end{tikzpicture}
\vspace{-0.8in}

\begin{center}
  {\color{white}\textbf{\LARGE AP Statistics}}\\[4pt]
  {\color{skymid}\large\textbf{Unit 2: Exploring Two-Variable Data}}\\[6pt]
  {\color{white}\textbf{\large Lesson 2.7 \quad Residuals}}
\end{center}

\vspace{0.22in}
\namedateperiod
\vspace{0.18in}

% ── LEARNING TARGETS ─────────────────────────────────────────────────────────
\begin{learningtargetbox}
\small
\begin{enumerate}[leftmargin=1.5em, itemsep=5pt]
  \item \ldots{} calculate a residual and interpret it in context as a measure of prediction
        error. \hfill\textcolor{slate}{\small(DAT-1.E)}
  \item \ldots{} construct a residual plot by graphing residuals vs.\ the explanatory variable
        (or predicted values). \hfill\textcolor{slate}{\small(DAT-1.E)}
  \item \ldots{} describe the pattern in a residual plot and use it to assess whether a linear
        model is an appropriate model for the data. \hfill\textcolor{slate}{\small(DAT-1.F)}
\end{enumerate}
\end{learningtargetbox}

\vspace{0.14in}

% ── TABLE OF CONTENTS ────────────────────────────────────────────────────────
\begin{tocbox}
\small
\renewcommand{\arraystretch}{1.5}
\begin{tabularx}{\linewidth}{@{} c l X r @{}}
\rowcolor{sky}
\textbf{\#} & \textbf{Component} & \textbf{Description} & \textbf{Score} \\
\midrule
1 & Warm-Up        & Spiral review --- regression equations and computing predicted values     & \blank{1.2cm} \\
2 & Guided Notes   & Residuals, residual plots, and assessing model appropriateness           & \blank{1.2cm} \\
3 & Group Activity & Computing residuals and interpreting residual plots (tiered)             & \blank{1.2cm} \\
4 & Exit Ticket    & Calculating and interpreting a residual; writing an AP-style conclusion  & \blank{1.2cm} \\
5 & Homework       & Practice computing residuals and using residual plots                    & \blank{1.2cm} \\
\midrule
  & \multicolumn{2}{r}{\textbf{Total}} & \blank{1.2cm} \\
\end{tabularx}
\end{tocbox}

\vspace{0.14in}

% ── KEEP IN MIND ─────────────────────────────────────────────────────────────
\begin{remindbox}
\small
A residual measures how far an actual data value falls from the regression line's prediction.
Three things the AP exam always asks you to do with residuals:

\vspace{6pt}
\begin{tabularx}{\linewidth}{@{} >{\centering\arraybackslash\columncolor{goldbg}}p{0.06\linewidth}
                                    X @{}}
\textbf{\large!} &
\textbf{Residual $=$ actual $-$ predicted ($y - \hat{y}$), \textit{not} $\hat{y} - y$.}\enspace
A positive residual means the actual value is above the regression line (the model
underpredicted). A negative residual means the actual value is below the line (overpredicted). \\[6pt]
\textbf{\large!} &
\textbf{Label your residual plot axes.}\enspace
The $x$-axis should show the explanatory variable (or $\hat{y}$); the $y$-axis should be
labeled ``residual.'' Draw a dashed horizontal reference line at $e = 0$. \\[6pt]
\textbf{\large!} &
\textbf{A complete residual-plot interpretation has three parts.}\enspace
(1) Describe the pattern (random scatter, curved, fanning). (2) State whether a linear model
\textit{is} or \textit{is not} appropriate. (3) Name the variables for context.
\end{tabularx}
\end{remindbox}

\end{document}
